% TOP_PROBLEM: {"id": 8400020, "title": "O18a — Recognizing primitive roots deterministically", "category_id": 3, "category": {"id": 3, "name": "graph_theory", "display_name": "Graph Theory", "description": "Problems involving graphs, networks, and their properties.", "slug": "graph-theory", "order_index": 3, "created_at": "2026-07-31T15:26:25.671Z"}, "status": "open", "problem_number": "AMR-083-0020", "set_id": 13}
% FIRST_POSTED: 2026-10-01
% ENTRY_KIND: statement_only
% UnsolvedMath ID 8400020; source catalogue code AMR-083-0020
% !TeX program = lualatex
% Definition and sources only; this file is not an LLM solution attempt.
\documentclass[11pt,a4paper]{article}
\usepackage[margin=25mm]{geometry}
\usepackage{fontspec}
\setmainfont{lmroman10-regular.otf}[BoldFont=lmroman10-bold.otf,
 ItalicFont=lmroman10-italic.otf,BoldItalicFont=lmroman10-bolditalic.otf]
\usepackage{amsmath,amssymb,mathtools}
\usepackage{unicode-math}
\setmathfont{latinmodern-math.otf}
\AtBeginDocument{\let\setminus\smallsetminus}
\usepackage{xurl,hyperref}
\hypersetup{colorlinks=true,urlcolor=blue}
\setcounter{secnumdepth}{0}
\setlength{\parindent}{0pt}
\setlength{\parskip}{5pt}
\setlength{\emergencystretch}{3em}
\begin{document}
\section*{O18a — Recognizing primitive roots deterministically}\label{8400020}
{\small UnsolvedMath ID 8400020 · AMR-083-0020 · Graph Theory · AMR Open Problem Lists}
\subsection{Definitions and mathematical statement}
For a prime $p$, the group $\mathbb F_p^\times$ consists of the nonzero residues modulo $p$ under multiplication. It is cyclic of order $p-1$. A residue $b$ is a primitive root if its powers exhaust this group, or equivalently if its multiplicative order is exactly $p-1$.

Given $p$ in binary and a residue $b\in\{0,\ldots,p-1\}$, can one decide deterministically, in time polynomial in $\log p$, whether $b$ is a primitive root? The answer for $b=0$ is negative; when $p=2$, the element $1$ generates the trivial unit group and is accepted. The primality of $p$ is an input promise.

An equivalent criterion for nonzero $b$ is
\[
b^{(p-1)/q}\not\equiv1\pmod p
\quad\text{for every prime }q\mid(p-1).
\]
This criterion specifies exactly the required property, but the distinct prime divisors of $p-1$ are not supplied as additional input. An algorithm that uses their factorization must include the cost of obtaining it.

The output is only a decision bit. Constructing a generator or computing the complete multiplicative order is not demanded. The proposed algorithm must have one unconditional polynomial bound covering every prime, including those for which $p-1$ has large prime factors. Random sampling, nonuniform precomputed tables, or successful verification of some specially chosen primes does not supply the universal deterministic decision procedure requested.
\subsection{Short English statement}
Can primitive roots modulo a given prime be recognized deterministically in polynomial bit complexity?
\subsection{Sources}
\begin{itemize}
\item[S1] ULAM.AI and the UnsolvedMath Contributors, supplied problems.json, record 8400020 (AMR-083-0020), AMR Open Problem Lists. Catalogue identity and source transcription; CC BY 4.0.
\url{https://huggingface.co/datasets/ulamai/UnsolvedMath}.
\item[S2] Adleman \& McCurley - Open Problems in Number Theory Complexity II (1994), O18a, p14 / LNCS p304. Original source indexed by AMR-083-0020.
\url{https://mccurley.org/papers/open.ps.gz}.
\end{itemize}
\end{document}
